\begin{proposition}[Crystalline local tangent condition]
\label{prop:selmer-crystalline-components}
Let $O$ be the integers of a finite extension $E/\mathbf Q_\ell$,
let $R$ be a complete Noetherian local quotient of
$O[[z_1,\ldots,z_d]]$, and let
$\rho:G_L\to\mathrm{GL}_3(R)$ be continuous, where
$L/\mathbf Q_\ell$ is finite.  Fix $e\ge1$, an ideal $J\subset R$,
an element $\delta\in R$, and an ideal
\[
                 I_t\supseteq\sqrt{J:\delta^\infty}.
\]
Suppose that every finite field-valued point of $R/J$ at which
$\delta\ne0$ gives a representation with a flag whose successive
characters $\kappa_0,\kappa_1,\kappa_2$ satisfy
\[
 \kappa_i|_{I_L}=\chi^{(1-i)e}|_{I_L},\qquad
 \kappa_0/\kappa_1\ne\chi,\qquad
 \kappa_1/\kappa_2\ne\chi,
\]
where $\chi$ is the cyclotomic character of $G_L$.
Then every point of $R/I_t$ over a finite $E$-algebra is crystalline
with Hodge--Tate weights $-e,0,e$ at every embedding of $L$.
In particular, its tangent cocycles satisfy the finite local condition.
\end{proposition}

\begin{proof}
For a crystalline representation $W$ with strictly negative
Hodge--Tate weights, local Tate duality and the Bloch--Kato dimension
formula give
\[
 \dim H^1(L,W)-\dim H_f^1(L,W)
               =\dim\operatorname{Hom}_{G_L}(W,E'(1)),
\]
where $E'$ is its finite coefficient field
\cite{bloch-kato-tamagawa,feng-local-bloch-kato}.
Every $\kappa_i$ is an unramified twist of a Tate character and is
crystalline.  Applying the formula to $\kappa_0/\kappa_1$ shows that
the first two steps $V_2$ of the flag form a crystalline extension.
The representation $V_2\otimes\kappa_2^{-1}$ then has successive
characters $\kappa_0/\kappa_2$ and $\kappa_1/\kappa_2$, of positive
Tate exponents $2e,e$.  Neither character is $\chi$: for the first,
their inertia exponents differ; the second is excluded by hypothesis.
Applying $\operatorname{Hom}_{G_L}(-,E'(1))$ to this flag gives zero,
even when the flag does not split.  Thus its entire first cohomology
is finite, and the last extension is crystalline as well.  Its Hodge
type is the same fixed multiset $\{-e,0,e\}$ at every point.

Put $I=J:\delta^\infty$ and $A=R[1/\ell]$.  Multiplication by
$\delta$ on $R/I$ is injective.  Therefore $\delta$ belongs to no
minimal prime: localization at such a prime is Artinian local,
where an element of the maximal ideal is nilpotent and cannot act
injectively.  Consequently $D(\delta)$ is dense in every component
of $V(IA)$.  The ring $A$ is Jacobson, with finite extensions of
$E$ as its closed residue fields
\cite{conrad-generic-deformation-fibres}.  Their maps on $R$ are
automatically local and continuous.  Hence the crystalline
field points considered above are dense in $V(IA)$.

Kisin's theorem supplies an ideal $K\subset A$ cutting out
crystallinity of this fixed Hodge type, with trivial inertial type
over $L$, for \emph{every finite $E$-algebra}
\cite{kisin-potentially-semistable}.
One first removes the $\ell$-power torsion from $R$ if necessary;
this leaves $A$ unchanged.  Since the dense field points kill $K$,
\[
                         K\subseteq\sqrt{IA}\subseteq I_tA.
\]
Every finite algebra-valued point of $R/I_t$ therefore factors
through Kisin's quotient.  This includes points over the dual
numbers, even above the locus $\delta=0$.

A crystalline dual-number representation is a crystalline
self-extension of its specialization.  Applying crystalline periods
splits its period sequence over $L_0\otimes_{\mathbf Q_\ell}E'$, a
product of fields.  Comparison then gives an $E'$-linear splitting
over $B_{\mathrm{cris}}$.  Its tangent class is therefore in $H_f^1$.
\end{proof}

The field-flag hypothesis follows directly from the ordinary equations.
Write $a_i(s)=\chi(s)^{(1-i)e}$ for $s\in I_L$, and choose
$t\in I_L$ such that $p_i=a_i(t)$ are distinct.  Those equations say
that $\rho(t)$ has characteristic roots $p_i$ and that
\[
  \prod_{i=0}^{2}\bigl(\rho(s_i)-a_i(s_i)1\bigr)=0
                   \qquad(s_i\in I_L).
\]
At a field point, diagonalize $D=\rho(t)$ in the order $p_0,p_1,p_2$.
Put $D_i=D-p_i1$ and $c_i=\prod_{j\ne i}(p_i-p_j)$.
For $A=\rho(s)$, the three substitutions $(s,t,t),(t,s,t),(t,t,s)$
give
\[
 (A-a_0(s)1)D_1D_2=0,\quad
 D_0(A-a_1(s)1)D_2=0,\quad
 D_0D_1(A-a_2(s)1)=0.
\]
The first fixes the first column, the last fixes the last row, and
the middle fixes $A_{11}$; hence $A$ is upper triangular with diagonal
$a_0(s),a_1(s),a_2(s)$.  Distinct eigenvalues make its common $a_0$
line and $a_2$ quotient plane unique.  Inertia is normal and the
$a_i$ are conjugation invariant, so the whole local group preserves
this flag.

For $\sigma\in G_L$, spectral projectors give
\[
 \begin{aligned}
 T_i&=\operatorname{tr}(D_{i+1}D_{i+2}\rho(\sigma))
       =c_i\kappa_i(\sigma),\\
 d_i&=c_{i+1}T_i-\chi(\sigma)c_iT_{i+1}\\
    &=c_ic_{i+1}\kappa_{i+1}(\sigma)
       \left(\frac{\kappa_i(\sigma)}{\kappa_{i+1}(\sigma)}
                         -\chi(\sigma)\right).
 \end{aligned}
\]
where indices on $D_i,c_i,T_i$ are cyclic and $i=0,1$ in the
second formula.  These are the factors defining
$\delta=\ell d_0d_1$ in the ordinary equations.
Thus $\delta\ne0$ supplies the two required inequalities.
The selected local component ideal contains
$\sqrt{J:\delta^\infty}$, so the proposition applies directly.
The restriction to $G_L=\ker\eta\cap\ker\tau$ removes the finite
inertial characters; its Hodge type remains the fixed parallel
multiset $\{-e,0,e\}$.
If $e\ge2$, the two nonresonance inequalities already follow from
inertia.  Taking $\delta=1$ then proves crystallinity on the reduced
ordinary locus itself; saturation is unnecessary for this local step.

The finite condition descends from this subgroup to the full local
group by averaging a $B_{\mathrm{cris}}$-valued primitive over its
finitely many cosets.  Here the restricted adjoint module is exactly
$\operatorname{End}(V)$, with its conjugation action.
Away from $\ell$, the same averaging makes
a tangent cocycle unramified whenever it vanishes on a subgroup of
finite index in inertia.  Thus the required local Selmer conditions
follow once the global deformation ring factors through these local
rings and imposes the stated finite inertia conditions.  These are
explicit conditions on the arithmetic model.
